\section{Native Implementation}

\subsection{Artifact boundary}

The runtime consumes the immutable VoiceDesign snapshot at revision
\code{5ecdb67327fd37bb2e042aab12ff7391903235d3}. Build-time and startup checks
validate configuration fields, tensor metadata, artifact hashes, and the
native ABI before serving traffic. The production image carries the
VoiceDesign weights and decoder-only speech-tokenizer material required for
inference. It excludes the tokenizer encoder, Base and CustomVoice checkpoints,
reference-audio functionality, and 0.6B weights.

\subsection{Talker and residual predictor}

The Rust model layer parses configuration and tokenizer artifacts, constructs
the VoiceDesign prompt from text, instruction, and official language name, and
owns the request lifecycle. Shared weights are uploaded once. Each session
retains independent key-value caches, workspaces, counters, random state, CUDA
stream, and cuBLAS handle.

For each codec frame, the talker predicts the semantic code and the native
15-step predictor produces the residual codebook values, yielding 16 codes per
frame. Sampling state is local to the request and is initialized from the
effective seed. Predictor and talker controls remain distinct; the predictor's
repetition penalty is fixed by the implemented model contract.

\subsection{Device-to-device handoff}

The production ABI avoids a codec-token round trip through host memory. The
talker exposes its device-resident frame and records a producer-ready CUDA
event. A non-blocking staging stream waits on that event, copies each frame into
a bounded device packet, and records completion. The decoder waits on the
packet-ready event and records a consumed event before the staging slot can be
reused. Tickets and device identity are checked so stale or cross-device state
cannot be accepted accidentally.

\subsection{Incremental decoder}

The decoder shares immutable weights but gives each active request its own
incremental neural state, stream, cuBLAS handle, events, histories, and bounded
PCM staging. CUDA kernels and cuBLAS execute the decoder path and clamp output
into signed 16-bit PCM. The current implementation uses non-blocking streams,
events, custom kernels, and cuBLAS. It does not claim CUDA Graph capture or
replay.

\subsection{Production image}

The runtime container targets \code{linux/arm64} and real \code{sm\_121} SASS
without a PTX fallback. The final stage contains CUDA runtime and the required
cuBLAS families but excludes compilers, development packages, Python, Node.js,
PyTorch, SGLang, vLLM, TensorRT, cuDNN, NPP, cuSPARSE, and NCCL. It runs as an
unprivileged user and is designed for a read-only root filesystem with dropped
Linux capabilities.
